% !TEX root = ../main.tex
\section{Marco teórico y ejercicios resueltos}

\subsection{Definición del problema y garantía analítica de factor 2}
Dado un grafo no dirigido $G = (V, E)$, una \textbf{cobertura de vértices} es un subconjunto $C \subseteq V$ tal que toda arista $(u, v) \in E$ satisface $u \in C \lor v \in C$. El problema \textit{Minimum Vertex Cover} busca minimizar la cardinalidad $|C|$ y es NP-hard.

El algoritmo voraz construye un matching maximal $M \subseteq E$ y añade ambos extremos de cada arista seleccionada a $C$. Dado que cada arista de $M$ aporta 2 vértices y cualquier cobertura válida requiere al menos un vértice por arista disjunta de $M$:
\begin{equation}
    |C| = 2|M| \le 2\,\text{OPT} \implies r(I) = \frac{|C|}{\text{OPT}} \le 2
\end{equation}

% -------------------------------------------------------------------------
\subsection{Ejercicio resuelto 1: validación de una cobertura de vértices}
El primer requisito consiste en verificar si un conjunto propuesto de vértices $C$ constituye una cobertura factible sobre el conjunto original de aristas del grafo.

\begin{lstlisting}[language=Python, caption={Código del Ejercicio Resuelto 1 (función de validación y pruebas)}, label={lst:es_vc}]
def es_vertex_cover(vertices, aristas, cobertura):
    cobertura = set(cobertura)
    if not cobertura.issubset(set(vertices)):
        return False
    for u, v in aristas:
        if u not in cobertura and v not in cobertura:
            return False
    return True

# Caso de prueba oficial de la guia: camino de 5 vertices
vertices = {"a", "b", "c", "d", "e"}
aristas = [("a", "b"), ("b", "c"), ("c", "d"), ("d", "e")]

print("Prueba 1 ({'b', 'd'}):", es_vertex_cover(vertices, aristas, {"b", "d"}))
print("Prueba 2 ({'a', 'c'}):", es_vertex_cover(vertices, aristas, {"a", "c"}))
\end{lstlisting}

\begin{lstlisting}[style=consolestyle, caption={Salida de consola del Ejercicio Resuelto 1}]
Prueba 1 ({'b', 'd'}): True
Prueba 2 ({'a', 'c'}): False
\end{lstlisting}

\subsubsection*{Análisis de resultados y validación}
\begin{itemize}
    \item El conjunto $\{b, d\}$ es una cobertura \textbf{válida} (\texttt{True}) porque el vértice $b$ cubre $(a, b)$ y $(b, c)$, mientras que $d$ cubre $(c, d)$ y $(d, e)$.
    \item El conjunto $\{a, c\}$ es \textbf{inválido} (\texttt{False}) porque la arista $(d, e)$ no tiene ninguno de sus extremos en el conjunto.
    \item \textbf{Principio de validación:} La comprobación debe efectuarse siempre sobre la lista de aristas originales completas sin descartar las aristas eliminadas durante la ejecución del algoritmo aproximado.
\end{itemize}

% -------------------------------------------------------------------------
\subsection{Ejercicio resuelto 2: algoritmo 2-aproximado y matching maximal}
El algoritmo recorre las aristas del grafo; cuando encuentra una arista cuyos dos extremos están libres, la incorpora al matching $M$ y añade ambos extremos a la cobertura $C$.

\begin{lstlisting}[language=Python, caption={Código del Ejercicio Resuelto 2 (algoritmo aproximado y ejecución)}, label={lst:vc_aprox}]
def vertex_cover_2_aprox(vertices, aristas):
    cobertura = set()
    matching = []
    for u, v in aristas:
        if u not in cobertura and v not in cobertura:
            matching.append((u, v))
            cobertura.add(u)
            cobertura.add(v)
    assert es_vertex_cover(vertices, aristas, cobertura)
    return cobertura, matching

# Ejecucion sobre el camino de 5 vertices
cobertura, matching = vertex_cover_2_aprox(vertices, aristas)
print("Matching maximal M:", matching)
print("Cobertura C:", sorted(list(cobertura)))
print("Tamanio |C|:", len(cobertura))
\end{lstlisting}

\begin{lstlisting}[style=consolestyle, caption={Salida de consola del Ejercicio Resuelto 2}]
Matching maximal M: [('a', 'b'), ('c', 'd')]
Cobertura C: ['a', 'b', 'c', 'd']
Tamanio |C|: 4
\end{lstlisting}

\subsubsection*{¿Por qué funciona la condición voraz?}
Una arista $(u, v)$ es seleccionada únicamente si ninguno de sus extremos pertenece a la cobertura en construcción. Como cada arista seleccionada incorpora inmediatamente sus dos extremos a $C$, cualquier arista posterior que comparta un vértice con las anteriores es ignorada. Por ende, las aristas almacenadas en $M$ son mutuamente disjuntas y forman estrictamente un \textbf{matching maximal}.

% -------------------------------------------------------------------------
\subsection{Ejercicio resuelto 3: búsqueda óptima exacta y razón de aproximación}
Para contrastar la calidad de la solución aproximada con la cota teórica, se implementa una búsqueda por enumeración exhaustiva de combinaciones para calcular el tamaño óptimo $\text{OPT} = |C^*|$.

\begin{lstlisting}[language=Python, caption={Código del Ejercicio Resuelto 3 (búsqueda óptima y cálculo de razón)}, label={lst:vc_exacto}]
from itertools import combinations

def vertex_cover_exacto(vertices, aristas):
    lista = sorted(list(vertices))
    for k in range(len(lista) + 1):
        for candidatos in combinations(lista, k):
            if es_vertex_cover(vertices, aristas, candidatos):
                return set(candidatos)
    return set()

# Comparacion experimental de calidad
optimo = vertex_cover_exacto(vertices, aristas)
aprox, matching = vertex_cover_2_aprox(vertices, aristas)
razon = len(aprox) / len(optimo) if optimo else 1.0

print("Optimo exacto C*:", sorted(list(optimo)), "con tamanio OPT =", len(optimo))
print("Solucion aproximada C:", sorted(list(aprox)), "con tamanio |C| =", len(aprox))
print(f"Razon observada r(I) = |C| / OPT: {razon:.2f}")

# Comprobacion formal de garantia
assert len(aprox) <= 2 * len(optimo), "Violacion de garantia teorica"
print("Garantia |C| <= 2*OPT: SATISFECHA RIGUROSAMENTE")
\end{lstlisting}

\begin{lstlisting}[style=consolestyle, caption={Salida de consola del Ejercicio Resuelto 3}]
Optimo exacto C*: ['b', 'd'] con tamanio OPT = 2
Solucion aproximada C: ['a', 'b', 'c', 'd'] con tamanio |C| = 4
Razon observada r(I) = |C| / OPT: 2.00
Garantia |C| <= 2*OPT: SATISFECHA RIGUROSAMENTE
\end{lstlisting}

\subsubsection*{Análisis de la razón observada y límites de complejidad}
\begin{itemize}
    \item Para el grafo camino $P_5$, el óptimo es $C^* = \{b, d\}$ con $|\text{OPT}| = 2$.
    \item El algoritmo aproximado produjo $|C| = 4$. Por tanto, la razón observada alcanza el peor caso teórico del algoritmo:
    \begin{equation}
        r(I) = \frac{|C|}{|\text{OPT}|} = \frac{4}{2} = 2.0
    \end{equation}
    \item \textbf{Costo computacional:} La búsqueda exhaustiva evalúa hasta $2^{|V|}$ subconjuntos ($O(2^{|V|} \cdot |E|)$). En contraste, el algoritmo 2-aproximado se ejecuta en tiempo estrictamente lineal $O(|E|)$, haciendo viable la resolución de instancias de gran escala donde el método exacto es intratable.
\end{itemize}
